\section{Introduction}

Debiasing methods like Just Train Twice (JTT) \cite{nam2020learning} succeed by reweighting examples that are learned later in training, implicitly relying on temporal ordering between spurious and core features. While JTT validates this operationally (reweighting works), the mechanistic foundation---when features converge at the gradient level---remains unmeasured. Without direct gradient-level verification, we cannot distinguish whether reweighting methods succeed due to temporal ordering or other factors (e.g., example hardness, loss landscape geometry), and cannot design principled interventions that exploit learning dynamics.

The problem extends beyond fairness metrics. Standard training (ERM) on spurious correlation benchmarks like Waterbirds \cite{sagawa2020distributionally} achieves 97\% accuracy on majority groups (waterbirds on water, landbirds on land) but collapses to 40\% on minority groups (waterbirds on land), where background cues mislead the model. Existing debiasing methods---GroupDRO \cite{sagawa2020distributionally}, Invariant Risk Minimization \cite{arjovsky2019invariant}, and JTT \cite{nam2020learning}---improve worst-group accuracy empirically, but lack mechanistic grounding in optimization dynamics. The temporal hypothesis---that spurious features converge earlier than core features during gradient descent---is stated in JTT and Learning from Failure \cite{liu2021just} papers but never verified via gradient-level measurement. Prior work tracked example-level hardness or forgetting \cite{toneva2018empirical}, not feature-level convergence.

The gap is both conceptual and methodological. Conceptually, we lack direct evidence that gradient descent's implicit bias toward simpler features \cite{soudry2018implicit} manifests as temporal separation between spurious and core feature convergence. Methodologically, measuring when features (not examples) converge requires isolating feature types---spurious-only vs core-only training variants---then tracking per-epoch gradient norms until stabilization. No prior work has performed this measurement.

We hypothesize that spurious features converge earlier not just heuristically, but measurably via per-epoch gradient norm tracking on ablated feature types. If spurious features (color, background) converge at epoch $E_s$ and core features (shape, semantics) converge at $E_c$ with $E_c - E_s \geq 2$ epochs, gradient descent's implicit bias toward simpler features creates an exploitable temporal gap. This gap arises from the feature complexity hierarchy: simpler low-level features (color, texture) stabilize in early convolutional layers before higher-level semantic features (shape, morphology) that require deeper hierarchical processing.

Building on this insight, we make the following contributions:

\textbf{1. Spurious features converge 4 epochs earlier than core features} ($\Delta=4$, substantially exceeding threshold by 2$\times$ in single-dataset proof-of-concept), providing the first direct gradient-level validation of the temporal ordering hypothesis underlying JTT/LfF. This mechanistic evidence explains why reweighting late-learned examples succeeds: it implicitly corrects for temporal misalignment driven by gradient descent's bias toward simpler features.

\textbf{2. Early convolutional layers drive temporal gaps} via significantly higher spurious correlation than late layers ($\rho_j=0.003$ vs $0.001$, $p=0.0028$, Cohen's $d=0.25$), confirming architectural feature hierarchy as mechanism. Layer-wise gradient analysis validates that temporal separation arises from hierarchical processing, not dataset artifacts.

\textbf{3. Cross-dataset transfer fails catastrophically} (39\% vs 86\% target worst-group accuracy), revealing that neuron correlations cannot be reused across spurious feature types---a fundamental constraint for gradient-aware methods. This negative result identifies the theoretical boundary: neuron-level interventions require per-dataset calibration, limiting scalability compared to example-level reweighting.

\textbf{4. Feature-level forgetting rate} provides independent stability metric (51\% lower for spurious features), confirming that spurious features converge to more stable predictions consistent with simpler decision boundaries.

Our work establishes that temporal ordering is measurable, mechanistically grounded in architectural feature hierarchy, and explains why reweighting methods (JTT, LfF) succeed: they implicitly correct for temporal misalignment by upweighting late-learned examples containing core features. However, neuron-level gradient modulation faces dataset-specificity constraints---$\rho_j$ values cannot be reused across spurious types, limiting the scalability of fine-grained interventions. The mechanistic validation opens future directions for adaptive reweighting schedules (monitor gradient convergence in real-time, adjust example weights dynamically) and architectural interventions (enforce late-layer learning delays to prevent early spurious dominance).

\textbf{Scope and limitations.} Results are validated on CMNIST color-based spurious correlation (single-dataset proof-of-concept with seed 0; full 10-seed statistical validation in progress). Generalization to background (Waterbirds), attribute (CelebA), and context (NICO++) spurious types is future work. Multi-dataset validation requires dataset-specific ablation strategies but uses identical gradient convergence measurement protocol.

We organize the paper as follows: Section 2 positions our work against prior debiasing methods and learning dynamics research. Section 3 describes our ablation training methodology and convergence measurement protocol. Section 4 details experimental setup across CMNIST and planned multi-dataset validation. Section 5 presents temporal gap results ($\Delta=4$ epochs), layer-wise mechanism validation, and intervention failure analysis. Section 6 discusses limitations (single-dataset scope, proof-of-concept statistical validation) and future multi-dataset generalization. Section 7 concludes by connecting back to the JTT/LfF temporal assumption and outlining adaptive intervention possibilities.
